\cotefile{157-5}
\scanpage{36}\mspage{23}

\unsure{suffisantes} : bien ordonné dans un cas, loc fini dans un autre. L'idée naturelle, c'est de \deleted{\unsure{prouver}} donner un énoncé commun, qui coiffe les deux à la fois, avec $I$ \unsure{artinien}, i.e. toute suite décroissante est stationnaire.

Ce qu'il faut, c'est une hyp. sur $I$ d'une part, sur $u_\bullet:X_\bullet\longrightarrow \unsure{Z|Y}_{\unsure{\bullet}}$ dans $M(I^{\unsure{\circ|o}})$ d'autre part, qui assure ceci

$(\unsure{\alpha|*})$\qquad $\forall\,i_0\in I$, et toute partie ouverte $J'$ de $J=I_{/i_0}$, la flèche $u_{J',J}$ dans le diagramme

\diagram{\begin{tikzcd}[ampersand replacement=\&,column sep=large]
 X_{J'} \arrow[r] \arrow[d,"u_{J'}"'] \arrow[dr,phantom,"\text{cocart}",pos=.54] \&
 X_J \arrow[d,dash,"\widetilde u_{J',J}"] \arrow[dd,bend left=45,"u_J"] \\
 Y_{J'} \arrow[r] \& Y_{J',J} \arrow[d,"u_{J',J}"'] \\
 \& Y_J
\end{tikzcd}}

est dans $C$.

Question : est-il exact qu'il suffit, pour la validité de \unsure{$\alpha$}, de supposer

\[
\begin{aligned}
&\text{a)}\quad u_\bullet\in \unsure{\mathbb C|C|\mathcal C}_I \quad\text{i.e.}\quad u_{J',J}\in C
 \quad \text{quand } J\deleted{\unsure{\in}}\unsure{=|\in}I_{/i_0},\\
&\qquad J'=J\setminus\{i_0\},\quad\text{et}
\end{aligned}
\]

b)\quad $I$ \uline{loc. fini} \deleted{\unsure{?}}

% archive sheet number: 35

\scanpage{37}\mspage{24}

\uline{Lemme} \deleted{\unsure{Pour|Soit} $i_0\in I$, \unsure{$K$} \unsure{désignons par} $P_{i_0}$} \inserted{Soit \deleted{\unsure{$I$}} $P\subset I$ l'ens des} $i_0\in I$ tels que pour tt ouvert $\deleted{J}K$ de $J=I_{/i_0}$, la flèche $u_{\deleted{\unsure{J'}}K,J}$ est dans $C$. \deleted{Pour qu'on ait} \inserted{Ceci dit,} soit $i_0\in\deleted{\unsure{P}}I$ tel que tout $i<i_0$ soit dans $P$. Alors $i_0\in P$. (On suppose $u_\bullet\in \unsure{C|\mathcal C}_I$, mais $I$ quelc.)

En effet, \deleted{\unsure{pour un tel cas}} soit $\deleted{J}K\subset J=I_{/i_0}$\inserted{ un ouvert}. Si \deleted{$\unsure{J'=J}$} $\deleted{\unsure{J'}}K=J$, il est trivial que $u_{\deleted{\unsure{J'}}K,J}\in C$ (c'est un iso). Si $\deleted{\unsure{J'}}K\ne J$, on a \deleted{$K\subset\unsure{J'}$}

\[
 K\subset J'\overset{\mathrm{def}}{=}\deleted{\unsure{I}}J\setminus\{i_0\}\underset{\unsure{\ne}}{\subset}J
\]

et pour prouver que $u_{K,J}\in C$, il suffit de prouver que $u_{K,J'}$, $u_{J',J}\in C$. Or $u_{\deleted{\unsure{K}}J',J\deleted{\unsure{q}}}\in C$ par hyp. $u_\bullet\in C_I$. Et $u_{\deleted{\unsure{J}}K,J'}\in C$ \unsure{par|pour} ce qui a été prouvé pour le cas de $J'$, compte tenu que \inserted{par hyp sur $i_0$,} tout $i\in J'$ (i.e. tt $i<i_0$) \deleted{\unsure{satisfait la}} \inserted{est dans $P$}.

\deleted{\unsure{Corollaire}} \uline{Corollaire}\qquad Si $I$ est artinien, $P=I$.

Car \unsure{toute} partie $P$ de $I$ ayant la propriété qu'on vient de dire ($I_{/i_0}'\subset P\Longrightarrow i_0\in P$) est égale à $I$, \deleted{\unsure{Car}} Si en effet il y avait $x_0\notin P$, il existerait $x_1<x_0$ avec $x_0\notin P$, et par récurrence on \unsure{construirait} une suite \unsure{strictmt|strictement} décroissante dans $I$, absurde.

% archive sheet number: 36

\scanpage{38}\mspage{25}

\S\,2. \uline{Description de} $C_I\subset\mathrm{Fl}(M(I^{\unsure{o|\circ}}))$, \uline{pour} $I$ \uline{ensemble ordonné}.

Soit $M$ une catégorie, stable par (petites) $\unsure{\underrightarrow{\lim}|\underline{\lim}}$, $C\subset\mathrm{Fl}\,M$ \inserted{(``cofibrations'').} On suppose

\[
\left\{\begin{aligned}
 &\text{a) } C\text{ contient les iso et stable par \unsure{cochgmt|cochangement} de base}\\
 &\text{b) } C\text{ stable par composition transfinie.}
\end{aligned}\right. \tag*{(2.1.)}
\]

Soit $I$ un \inserted{(petit)} ensemble ordonné, et considérons

\[
 M(I^o)=\underline{\mathrm{Hom}}\,(I;\deleted{\unsure{C}}M),
\]

on va définir une partie

\[
 C_I\subset\mathrm{Fl}\,(M(I^o))
\]

ainsi, pour $u_\bullet:X_\bullet\longrightarrow Y$ flèche \unsure{de|dans} $M(I^o)$

\[
\begin{aligned}
 u_\bullet\in C_I\ &\Longleftrightarrow\ \forall\,i_0\in I\ \text{et pour tt ouvert}\\
 &J\subset I_{/i_0}=K,\ \text{la flèche}\\
 &u_{J,\deleted{\unsure{I_{/i_0}}}K}\ \text{est}\ \in C
\end{aligned}\tag*{(2.2.)}
\]

Voici la définition de $u_{J,K}$. Pour tt ouvert $J$ de $I$, on pose

\[
\left\{\begin{aligned}
 X_J&=\underset{i\in J}{\underrightarrow{\lim}}\,X_{\unsure{i|\bullet}},\quad
 Y_J=\underset{i\in J}{\underrightarrow{\lim}}\,Y_i,\\[-2pt]
 &\deleted{\text{et on considère la}}\\[-2pt]
 u_J&:X_J\longrightarrow \deleted{\unsure{X}}Y_J,\quad
 u_J=\underset{i\in J}{\underrightarrow{\lim}}\,u_i,
\end{aligned}\right.\tag*{(2.3.)}
\]

et on considère, pour $\deleted{\unsure{J}}K\subset\deleted{\unsure{K}}J$, le diagramme

% archive sheet number: 37

\scanpage{39}\mspage{26}

\diagram{\begin{tikzcd}[ampersand replacement=\&,column sep=large,row sep=large]
\text{(2.4.)}\&
 X_{\deleted{J}\unsure{K|J}}\arrow[r,"\tau^X_{J,K}"] \arrow[d,"u_{\deleted{J}K}"'] \arrow[dr,phantom,"\text{cocart}",pos=.54]
 \& X_{\deleted{K}J}\arrow[d,"\widetilde u_{\deleted{J}K,\deleted{K}J}"]\arrow[dd,bend left=64,"u_{\deleted{K}J}"]\\
\& Y_{\deleted{J}K}\arrow[r,"\tau^Y_{J,K}"]
 \& Y^{\deleted{\unsure{u}}}_{\deleted{J}K,\deleted{K}J}=Y_{\deleted{J}K}\amalg_{X_{\deleted{J}K}}X_{\deleted{K}J}
 \arrow[d,"u_{\deleted{J}K,\deleted{K}J}"]\\
\& \&Y_{\deleted{K}J}
\end{tikzcd}}

\uline{Théorème} 2.1. Soit $M$ \inserted{$\unsure{\mathcal U|\mathscr U}$-}catégorie stable par petites $\unsure{\underrightarrow{\lim}|\underline{\lim}}$, $C\subset\mathrm{Fl}(M)$ satisfaisant les conditions de stabilité (2.1.) (par \unsure{cochgt} de base et par composition transfinie), $I$ un \inserted{petit} ensemble ordonné, $u_\bullet:X_\bullet\longrightarrow Y_\bullet$ une flèche dans $M(I^{\unsure{o|\circ}})$. \deleted{On}

$1^\circ$) \inserted{On suppose $u_\bullet\in C_I$ (cf. (2.2.)).} Soient $\deleted{J}K\subset\deleted{K}J$ deux ouverts de $I$. Alors $u_{\deleted{J}K,\deleted{K}J}$ (défini \unsure{ds|dans} (2.4.)) est dans $C$.

$2^\circ$) \inserted{Supposons $I$ artinien.} Pour qu'on ait $u_\bullet\in C_I$, i.e. pour que $u_{\deleted{J}K,\deleted{K}J}$ soit dans $C$ pour tt couple $(\deleted{J}K\subset\deleted{K}J)$ d'ouverts tels que $\deleted{K}J$ ait un plus grand élément \inserted{i.e. soit de la forme $I_{/i_0}$} (donc aussi, par $1^\circ$, pour \uline{tout} couple d'ouverts $J\subset K$), il (faut et) suffit qu'il en soit ainsi dans le cas particulier où $J=I_{/i_0}$, $J'=I'_{/i_0}$

\[
 \overset{\mathrm{def}}{=} I_{/i_0}-\{i_0\}=\{i\in I\mid i<i_0\}.
\]

% archive sheet number: 38

\scanpage{40}\mspage{27}

$3^\circ$) \inserted{Supposons encore $I$ artinien\textsuperscript{*}.} Soit $W\unsure{\subset|\in}\mathrm{Fl}(M)$ un localiseur modérément \deleted{\unsure{satisfait}}\unsure{saturé|saturisé}, tel que \deleted{\unsure{la}} $TC\overset{\mathrm{def}}{=}C\cap W$ satisfasse également aux conditions de stabilité (1.1.) (par \unsure{cochgt} de base, par composition transfinie).

% slanted
\margin{J'ignore si cette hypothèse est nécessaire, cf commentaire p. \unsure{39|35}.}

Les conditions suivantes sur une flèche $u_\bullet:\deleted{\unsure{Y}}X_\bullet\longrightarrow\unsure{Y|Z}$ dans $M(I^{\unsure{o|\circ}})$ sont équivalentes

a) $u_\bullet\in(TC)_I$, i.e. pour tout $i_0\in I$ et $K$ ouvert de $J=I_{/i_0}$, la flèche $u_{K,J}$ est dans $TC$

b) $u_\bullet\in C_I\cap W_{I^o}$, i.e. \deleted{Pour les} pour tout $i_0\in I$, on a $u_{i_0}\in W$, et pour tout ouvert $K$ de $I_{/i_0}=J$, on a $u_{K,J}\in C$.

Donnons tout de suite les corollaires principaux :

\uline{Corollaire \unsure{2.2.|2.1.}}: Supposons $(u_\bullet:X_\bullet\longrightarrow Y)\in C_I$. Alors

a) $\forall\,i\in I$, $u_i:X_i\longrightarrow Y_i$ est dans $C$.

b) \deleted{\unsure{$\forall$}} La flèche

\[
 u_I:X_I\longrightarrow Y_I\tag*{(2.2.1.)}
\]

est dans $C$. Plus généralement, pour \unsure{tte} partie ouverte $J$ de $I$, la flèche

\[
 u_J:X_J\longrightarrow Y_J\tag*{(2.2.2)}
\]

est dans $C$.

% archive sheet number: 39

\scanpage{41}\mspage{28}

Pour a), on applique $1^\circ$) \inserted{(ou simplement la \uline{déf.} de $C_I$)} dans le cas de l'inclusion d'ouverts $\varnothing\hookrightarrow I_{/i_{\unsure{0}}}$, \unsure{vu} qu'on a

\[
 u_{\varnothing,I_{/i_0}}\simeq u_{i_0}.
\]

Pour b), on applique $1^\circ$) à l'inclusion $\varnothing\longrightarrow I$, vu que

\[
 u_{\varnothing,I}\simeq u_I,
\]

et plus généralement

\[
 u_{\varnothing,J}\simeq u_J.
\]

\uline{Corollaire} 2.3. Supposons $u_\bullet\in C_I\cap W_{I^{\unsure{o|\circ}}}$. Alors \deleted{on} pour \unsure{tte} \deleted{\unsure{$i\in I$},} inclusion d'ouverts $K\subset J$ de $I$, on a $u_{\deleted{J}K,J}\deleted{\unsure{\text{est}}}\in TC$, et en particulier \deleted{(\unsure{3|2}.3.1)} (prenant $K=\varnothing$) on a $u_J\in TC$. Enfin pour tout $i\in I$, on a $u_i\in TC$.

\uline{Corollaire} 2.4. \deleted{\unsure{Supposons}} Soit $X_\bullet$ \unsure{de|dans} $M(I^o)_C$ un syst. inductif $C$-cofibrant, i.e. tel que $\varnothing_I\longrightarrow X_\bullet$ soit dans $C_I$, i.e. tel que \deleted{\unsure{les}} pour tout $i_0\in I$ et tout ouvert $K$ de $J=I_{/i_0}$, la flèche de transition $\tau_{K,J}:X_K=\underset{i\in K}{\underrightarrow{\lim}}\,X_i\longrightarrow X_J\simeq X_{i_0}$ soit dans $C$. Alors pour \uline{tout} couple d'ouverts $\deleted{J}K\subset\deleted{K}J$, la flèche \deleted{$\tau_{K,J}:$}

\[
 \tau_{K,J}:\underset{i\in K}{\unsure{\underrightarrow{\lim}|\underline{\lim}}}\,X_i
 \longrightarrow\underset{\unsure{i|i'}\in J}{\underrightarrow{\lim}}\,X_{\unsure{i'|i}}
\]

est dans $C$\textsuperscript{(*)}. \margin{et réciproquement, cela caractérise les objets $C$-cofibrants}
% slanted
En particulier (faisant $K=I_{/k}$, $J=I_{/j}$ pour $\unsure{h|k}\le j$) si $k\le j$ sont dans $I$, le morphisme de transition

\[
 \tau_{\unsure{k,j}}\quad X_k\unsure{\longrightarrow|-}X_j
\]

% archive sheet number: 40

\scanpage{42}\mspage{29}

est dans $C$. Enfin, si $W$\inserted{$\subset\mathrm{Fl}\,M$} est comme dans 2.1.$3^\circ$, et si, $\varnothing_I\longrightarrow X$ est dans $W_{I^{\unsure{o|\circ}}}$ i.e. les $X_i$ sont $\unsure{W|\acute W|\widetilde W}$-coasphériques, alors \deleted{\unsure{$\varnothing$}} $X_\bullet$ est dans $M(I^o)_{TC}$ (où $TC=C\cap W$), et pour tt ouvert $J$ de $I$, $X_{\deleted{I}J}=\underset{i\in\deleted{I}J}{\unsure{\underrightarrow{\lim}|\underline{\lim}}}\,X_i$ est dans $\inserted{M}_{TC}$, \inserted{donc $\unsure{X_j|X_J}$ est $W$-coasphérique,} et pour tt couple d'ouverts $K\subset J$ la flèche de transition $\tau_{K,J}$ est dans $TC$.

Signalons enfin le complément suivant au théorème 2.1., qui résultera de sa démonstration :

\uline{Corollaire} 2.5. Supposons que l'espace topologique associé à $I$ (\inserted{dont} les ouverts sont les ouverts de $I$, et les $I_{/i}$ forment une base d'ouverts) soit \uline{noethérien}, i.e. toute suite croissante d'ouverts est stationnaire. Alors \unsure{toute|toutes} les conclusions dans le th. 2.1. et ses corollaires restent \unsure{valables}, en supposant seulement la stabilité de $C$ ou de $TC=C\cap W$ par composition simple (au lieu de composition transfinie), en plus des deux autres stabilités.

% archive sheet number: 41

\scanpage{43}\mspage{30}

\margin{\unsure{$2^{\mathrm{e}}$|$\S$}}

Démonstration du th. 2.1 et du cor. 2.5.

2.1. $1^\circ$). \deleted{On prouve d'abord le cas particulier pour $K=\varnothing$, i.e. \unsure{que|pour} $\forall$ ouvert $J$ de $\deleted{I}\unsure{I|A}$, la flèche}

\deleted{\[
 u_{\unsure{J|I}}:X_J\longrightarrow Y_J
\]}

\deleted{est dans $C$. Considérons l'ensemble \inserted{$P\subset\unsure{\mathcal P|\mathfrak P}(I)$} des ouverts $J$ de $I$ pour lesquels cette condition soit satisfaite, prouvons que \inserted{pour} toute partie $(J_\alpha)_{\alpha\in A}$ tot. ordonnée \inserted{(par inclusion)} de $P$, la réunion $\bigcup J_\alpha$ est encore dans $P$. On sait (\unsure{où|ou} on \unsure{pourra|prouve} en fait \unsure{comment}) qu'il existe, pour tt ens. tot. ordonné, \deleted{de toute partie} \inserted{une partie} bien ordonnée cofinale. Donc on peut \unsure{suppo-} $A$ bien ordonné.}

\unsure{OPS} $J=I$, \inserted{$K\subsetneqq I$.} On construit une suite transfinie \inserted{str. croissante} d'ouverts $I_\alpha$, avec

\[
 \left\{\begin{aligned}
 I_0&=\deleted{\varnothing}K\\
 I_{\alpha+1}&= I_\alpha\cup I_{/s_\alpha}\quad\text{où }s_\alpha\in I\setminus I_\alpha
 \quad(\text{si }I_\alpha\ne I,\\[-2pt]
 &&\text{sinon la \unsure{constr.}}\\[-2pt]
 &&\text{transfinie s'arrête})\\
 I_\alpha&=\bigcup_{\beta<\alpha}I_\beta\quad\text{si $\alpha$ ordinal limite.}
 \end{aligned}\right.
\]

Il existe un $\alpha_0$ tel qu'on ait $I_{\alpha_0}=I$. On voit alors comme au \S\,1

% archive sheet number: 42

\scanpage{44}\mspage{31}

que pour tt $\alpha$, \deleted{$\unsure{u_{I_\alpha}}$} on a

\[
 u_{I_\alpha,I_{\alpha+1}}\in C,\tag*{$(\unsure{*})$}
\]

car il est déduit de $u_{I_{/s_\alpha}\cap I_\alpha,\,I_{/s_\alpha}}$ par cochangement de base (cf diagramme 6 p. 17, qu'il faudrait \unsure{dégager} \unsure{en|un} lemme, débarrassé des particularités techniques de la situation, i.e. un lemme sur les hom de carrés \unsure{cartésiens} :

\uline{Lemme} 2.6. Soit un hom de carrés cocartésiens \inserted{$X_\bullet,Y_\bullet$}, dans une catégorie $M$ \unsure{à $\cup$|où} \deleted{\unsure{lim}} sommes amalgamées (NB il suffit que les sommes amalgamées intervenant ci-dessous existent) :

\diagram{\begin{tikzcd}[ampersand replacement=\&,row sep=small,column sep=small]
\& X_0\arrow[rr]\arrow[dr]\arrow[dd,"f_0"'] \& \& X_1\arrow[dr]\arrow[dd,"f_1",near end] \&\\
\& \& X'_0\arrow[rr,crossing over]\arrow[dd,"f'_0"',crossing over] \& \& X'_1\arrow[dd,"f'_1"]\\
\text{(2.6.1)}\&Y_0\arrow[rr]\arrow[dr] \& \&Y_1\arrow[dr] \&\\
\& \&Y'_0\arrow[rr]\&\&Y'_1
\end{tikzcd}}

Soient $Z_0\unsure{\ne|,|;|\text{et}}Z_1$ les deux carrés latéraux et considérons, pour tt carré commutatif

\diagram{\begin{tikzcd}[ampersand replacement=\&,column sep=small,row sep=small]
Z={} \& X\arrow[r]\arrow[d] \& X'\arrow[d]\\
\&Y\arrow[r,dash]\&Y' \&\deleted{X}\arrow[r,dash]\&\deleted{\unsure{X'}}
\end{tikzcd}}

, la flèche $u_Z$ du diagramme

\diagram{\begin{tikzcd}[ampersand replacement=\&,column sep=small,row sep=small]
 X\arrow[r]\arrow[d]\arrow[dr,phantom,"\text{cocart}"] \& X'\arrow[d]\&\\
 \deleted{\unsure{X}}Y\arrow[r] \&\deleted{\unsure{Y}}\widetilde Z\arrow[r,"u_Z"']\&Y'\quad\unsure{\longrightarrow|\text{---}}\,\unsure{.}
\end{tikzcd}}

% archive sheet number: 43

\scanpage{45}\mspage{32}

Alors la flèche $u_{Z_1}$ se déduit de $u_{Z_0}$ par cochangement de base, de façon plus précise on a \deleted{\unsure{On}}un carré cocartésien, où

\diagram{\begin{tikzcd}[ampersand replacement=\&,column sep=small,row sep=small]
T={}\&X_0\arrow[r]\arrow[d]\&X_1\arrow[d]\\
\&Y_0\arrow[r,dash]\&Y_1
\end{tikzcd}}

et

\diagram{\begin{tikzcd}[ampersand replacement=\&,column sep=small,row sep=small]
T'={}\&X'_0\arrow[r,dash]\arrow[d]\&X'_1\arrow[d]\\
\&Y'_0\arrow[r]\&Y'_1
\end{tikzcd}}

désignent les faces \unsure{postérieure|postérieures} et \unsure{antérieure|antérieures} du cube (2.6.1) :

\diagram{\begin{tikzcd}[ampersand replacement=\&]
\&\widetilde Z_0\arrow[r,"u_{Z_0}"]\arrow[d]\arrow[dr,phantom,"\text{cocart}"]\&Y'_0\arrow[d]\\
(\unsure{*|\alpha})\&\widetilde Z_1\arrow[r,"u_{Z_1}"]\&Y'_1,
\end{tikzcd}}

\deleted{déduit du diagramme \unsure{multicocartésien}} qu'on voit sur la base du diagramme

\diagram{\begin{tikzcd}[ampersand replacement=\&,row sep=small,column sep=small]
\&X_0\arrow[rr]\arrow[dr,dash]\arrow[dd,dash,"f_0"']\&\&X_1\arrow[dr,dash]\arrow[dd,"f_1",near end]\&\&\\
\&\&X'_0\arrow[rr,dash,crossing over]\arrow[dd,dash]\arrow[dddr,"f^{\unsure{\prime}}_0",near start]\&\&X'_1\arrow[dd]\arrow[dddr,"f'_1"]\&\\
(\unsure{**|\alpha\alpha})\&Y_0\arrow[rr]\arrow[dr]\&\&Y_1\arrow[dr]\&\&\\
\&\&\widetilde Z_0\arrow[rr,dash]\arrow[dr,shift left=.5]\arrow[dr,shift right=.5,"u_{Z_0}"']\arrow[ur,phantom,"\text{cocart}",sloped]\&\&\widetilde Z_1\arrow[dr,"u_{Z_1}"']\&\\
\&\&\&Y'_0\arrow[rr]\arrow[ur,phantom,"\text{cocart}",sloped]\&\&Y'_1,
\end{tikzcd}}

dans lequel le cube est un cube totalement cocartésien, déduit du carré cocartésien $X_\bullet$ en le ``\unsure{copoussant}'' par la flèche $f_0:X_0\longrightarrow Y_0$. La base du cube, et le carré

% archive sheet number: 44

\scanpage{46}\mspage{33}

composé de celui-ci et de $(\unsure{*|\alpha})$ \deleted{\unsure{étant}} (base du diagramme \unsure{polyédral|polyèdral} $(**)$) étant \unsure{cocartésiens|cocartésien}, il en est de $=$ de $(*)$.

\uline{Lemme} 2.7. Soit $(I_\alpha)_{\alpha\le\alpha_0}$ une suite ordinale croissante d'ouverts de $I$, \deleted{dont réunion $I_{\alpha_0}$. Alors, \unsure{et|de}} Alors $u_{I_0,I_{\alpha_0}}$$:$\inserted{$Y_{I_0,I_{\alpha_0}}\unsure{\mathbin{-}|\longrightarrow}Y_{I_{\alpha_0}}$} est composée transfinie, indexée par l'ens des ordinaux $\deleted{\unsure{\alpha<\alpha_0}}\inserted{-1\le}\alpha<\alpha_0$, d'une suite de flèches $v_\alpha$, où chaque $v_\alpha$ est déduit de $u_{I_\alpha,I_{\alpha+1}}$ par cochangement de base (et de $u_{I_0}:X_{I_0}\longrightarrow Y_{I_0}$ pour $\alpha=0$).

Pour le voir, on introduit la suite transfinie d'objets $Z_\alpha$\inserted{$=Y_{I_\alpha,I_{\alpha_0}}$}, définis par le \deleted{$Z_\alpha=$} diagramme

\diagram{\begin{tikzcd}[ampersand replacement=\&,row sep=large,column sep=large]
 \deleted{\unsure{Z|X}}X_{I_\alpha}\arrow[r]\arrow[d] \& X_{I_{\alpha_0}}\arrow[d,dash]\\
 Y_{I_\alpha}\arrow[r] \&\deleted{\unsure{Z|Y}}\unsure{Y|Z}_{I_\alpha,I_{\alpha_0}}=Z_\alpha\arrow[d]\\
 \&Y_{I_{\alpha_0}},
\end{tikzcd}}

on a \deleted{\unsure{du}}

\[
 Z_0=Y_{I_0,I_{\alpha_0}},\quad Z_{\alpha_0}=Y_{I_{\deleted{0}\alpha_0},I_{\alpha_0}},\quad
 Z_\alpha=\underset{\beta<\alpha}{\unsure{\underrightarrow{\lim}|\underline{\lim}}}\,Z_\beta
 \quad\text{si }\unsure{\alpha}
\]

% archive sheet number: 45

\scanpage{47}\mspage{34}

ordinal limite. Donc il reste à déterminer les $\unsure{v|u}_{\alpha,\alpha+1}:Z_\alpha\longrightarrow Z_{\alpha+1}$\inserted{ (pour $\alpha<\alpha_0$)}, ce qui se voit sur le diagramme

\diagram{\begin{tikzcd}[ampersand replacement=\&,row sep=large,column sep=large]
 X_{I_\alpha}\arrow[r]\arrow[d] \& X_{I_{\alpha+1}}\arrow[r]\arrow[d] \& X_{I_{\alpha_0}}=Z_0\arrow[d,Rightarrow]\\
 Y_{I_\alpha}\arrow[r,dash] \& Y_{I_\alpha,I_{\alpha+1}}\arrow[r,dash]\arrow[d,"u_{I_\alpha,I_{\alpha+1}}"] \& \deleted{\unsure{Z|Y}}\unsure{Y|Z}_{I_\alpha,I_{\alpha_0}}=Z_\alpha\arrow[d,"v_{\alpha,\alpha+1}"]\\
 \& Y_{I_{\alpha+1}}\arrow[r] \& Y_{I_{\alpha+1},I_{\alpha_0}}=Z_{\alpha+1}.
\end{tikzcd}}

\deleted{\unsure{Donc} Ainsi, \unsure{Rest. Rev.}}

Reprenons la démonstration de 2.1 $\deleted{\unsure{1^\circ|2^\circ}}\unsure{2^\circ|1^\circ}$.

\deleted{Dans le cas actuel $I_0=\varnothing_M$, donc $X_{I_0}\longrightarrow Y_{I_0}$ est $\varnothing_{\unsure{M|0}}\xrightarrow{\sim}\varnothing_M$, donc $Z_0={}$}

On \deleted{\unsure{l'}}a prouvé, par les \inserted{2} lemmes précédents, \deleted{dans le cas particulier $K=\varnothing$, \unsure{On} procède de $=$ dans le cas général\unsure{,} d'un couple $(K,\unsure{J|I})$ quelconque, en utilisant que l'on sait déjà que
\[
 (u_K:X_K\longrightarrow Y_K)\in C.
\]}

\deleted{\unsure{Ceci} fait, en revenant \unsure{au} $I_0=K$, dans $I$.}

Notons que lorsque l'espace sous-jacent à $I$ est noethérien, alors la construction transfinie est forcément finie, donc \unsure{on} n'a pas besoin de la stabilité de $C$ par \unsure{composés|composées} transfinis infinis,

% archive sheet number: 46

\scanpage{48}\mspage{35}

Prouvons $2^\circ$. On suppose donc que pour \unsure{tt} $i_0\in I$, on ait $u_{I_{<i_0},I_{/i_0}}\in C$, on veut en déduire que $u_\bullet\in\deleted{\unsure{I}}C_I$ i.e. $u_{J,I_{/i_0}}\in C$ pour \uline{tout} ouvert. \deleted{Soit $P$ l'ens \unsure{d'indices} \unsure{qui sont} \unsure{des}}

\uline{Lemme} 2.8. (Sans hypothèse artinienne sur $I$) Soit $P\subset I$ l'ens des $i_0\in I$ tels que pour \deleted{\unsure{les}} $u_\bullet\mid I_{/i_0}$ soit dans $C_{I_{/i_0}}$, i.e. que pour tout $i\le i_0$, et tout ouvert $K$ de $I_{/i}$, on ait $u_{K,I_{/i}}\in C$. \deleted{Alors} Soit $i_0\in I$ tel que $I_{<i_0}\subset P$. Alors $i_0\in P$.

\uline{Dém} Il reste à prouver que si $\deleted{\unsure{J|I}}K\subset I_{/i_0}=J$ est un ouvert, alors $u_{K,\deleted{\unsure{I_{/i_0}}}J}\in C$. OPS $K\ne I_{/i_0}$, i.e. $K\subset I_{<\unsure{i_0|i'_0}}\overset{\mathrm{def}}{=}J'$. Mais alors par hypothèse initiale sur $u$, \inserted{\unsure{on a}} $u_{J',J}\in C$, et d'autre part, \deleted{\unsure{le}} par hyp.\inserted{ sur $i_0$, \deleted{\unsure{$J$|je}}} $u_\bullet\mid J'$ est dans $C_{J'}$. Donc par le Th.\ 2.1, $1^\circ$, on a $u_{K,J'}\in C$. Donc par transitivité $u_{K,J}\in C$, qed.

\uline{Corollaire} 2.\deleted{\unsure{8|9}}\unsure{9|8}. Si $I$ est artinien, alors $P=I$, i.e. $u_\bullet\in C_I$.

Cela prouve donc 2.1. $2^\circ$. Notons que $\operatorname{Esp}(I)$ est noethérien, il en est de même de \uline{tout} ouvert (tel $J'$ ci-dessus), donc

% archive sheet number: 47

\scanpage{49}\mspage{36}

pour la conclusion du lemme, inutile de supposer $C$ stable par compositions infinies.

Prouvons 2.1. $3^\circ$. On a $a\Longrightarrow b$, car par 2.1. $1^\circ$\unsure{,|)} $u_\bullet\in(TC)_I$ implique que pour $i\in I$, $u_i:X_i\longrightarrow Y_i$ est dans $TC$, \deleted{donc dans} (cf cor.\ 2.1. \unsure{a}), \unsure{Donc} il reste à prouver $\mathrm{b)}\Longrightarrow\mathrm{a)}$ i.e. si $u_\bullet\in C\cap W_{I^{\unsure{o|\circ}}}$, alors $u_\bullet\in(TC)_I$, i.e. les $u_{K,I_{/i_0}}$ sont dans $TC$ \deleted{\unsure{puisq.}}

On sait déjà que $u_{K,I_{/i_0}}\in C$ par hyp. $u_\bullet\in C_I$, donc il faut prouver que c'est $\in W$. On est ramené à prouver que $u_K\in W$, i.e. $u_K\in TC$, car \deleted{par} le diagramme (où $J=I_{/i_0}$)

\diagram{\begin{tikzcd}[ampersand replacement=\&,column sep=huge]
 X_K\arrow[r]\arrow[d,"u_K"']\arrow[dr,phantom,"\text{cocart}"] \& X_J\arrow[d,"\widetilde u_{K,J}"]\arrow[dd,bend left=60,"u_J"]\\
 Y_{K.}\arrow[r,dash] \& Y_{K,J}\arrow[d,"u_{K,J}"]\\
 \&Y_J
\end{tikzcd}}

$u_K\in TC$ implique $\widetilde u_{K,J}\in TC$, et comme $u_J\simeq u_{i_0}\in W$ par hypothèse, il en résulte bien $u_{K\unsure{J|,J}}\in W$ par saturation. \inserted{Mais on va}

\deleted{Donc on est ramené à prouver le \uline{Lemme} 2.9. \inserted{Supposons $I$ artinien.} Soit $u_\bullet\in C_I\cap W_{I^o}$. Alors pour \unsure{tt} ouvert $K$ de $I$, $u_K:X_K\longrightarrow Y_K$ est dans $\deleted{\unsure{T}}TC$ (donc dans $W$),} s'y prendre autrement:

% archive sheet number: 48

\scanpage{50}\mspage{3\deleted{\unsure{8|6}}7}

\uline{Lemme} 2.9. $I$ ens. ordonné quelconque, $C,W$ comme dans Th.\ 2.1, $1^\circ$ et $3^\circ$, $\deleted{\unsure{v}}u_\bullet:X_\bullet\longrightarrow Y_\bullet$ dans $C_I\cap W_{I^{\unsure{o|\circ}}}$. Soit $P$ l'ens des $i_0\in I$ tels que $u_\bullet\mid I_{/i_0}$ soit dans $(TC)_{I_{/i_0}}$. \deleted{Alors} Si $i_0\in I$ est tel que $I_{<i_0}\unsure{\in|\subset}P$, alors $i_0\in P$.

En effet, \deleted{\unsure{on aura} $u_\bullet\mid$} il reste à prouver que pour \uline{tout} ouvert $K$ de $J=I_{/i_0}$, \inserted{$K\ne J$,} on a $u_{K,J}\in TC$. Mais \deleted{\unsure{on a} $K\subset J'=I_{<i_0}$, et il suffit de voir que $u_{K,J'}\in TC$, $u_{J',J}\in W$.} Pour ceci, il suffit de voir que $u_K,u_J\in\deleted{\unsure{TC|W}}\unsure{W|TC}$ (comme on vient de le voir). Or $u_J\simeq u_{i_0}\in W$ par hypothèse.

D'autre part $\deleted{\unsure{u_K\in}}u_\bullet\mid K\in(TC)_K$ \deleted{par hyp. sur} puisque $K\subset I_{<i_0}\subset P$, donc par le Th.\ 1, $1^\circ$, on conclut que $u_K\in TC$.

Si maintenant on suppose $I$ artinien, il résulte de 3.9. que $P=I$, donc $u_\bullet\in TC_I$, qed.

Cela achève la démonstration du Th.\ 2.1, et aussi du cor.\ 2.5, vu que pour la \inserted{\unsure{dém.}} \inserted{du} lemme 2.9. également, les hypothèses infinies sur $C$, $C\cap W$ sont inutiles si $\operatorname{Top}(I)$ est noethérien, donc ses ouverts \deleted{\unsure{tt}} \deleted{\unsure{étant}} sont \unsure{aussi} noethériens.

\uline{Remarque}\deleted{\unsure{s}} $1)$ Dans 2.1 $2^\circ$, l'hypothèse $I$ artinien n'est pas superflue, comme on voit en prenant $I=\mathbb N^{\mathrm{opp}}\simeq\deleted{\unsure{\mathbb Z}}\mathbb Z^{\le0}$,

% archive sheet number: 49

\scanpage{51}\mspage{3\deleted{\unsure{9|8}}8}

et pour une flèche

\[
 \varnothing_{\unsure{M(I)}}\longrightarrow Y_\bullet
\]

Dire que cette flèche est dans $C_I$ signifie que les morphismes de transition $\deleted{\unsure{x}}Y_i\xrightarrow{u_i}Y_{i+1}$ sont dans $\deleted{\unsure{T}}C$, \uline{et} que les $Y_i$ sont $C$-cofibrants i.e. $\varnothing_M\longrightarrow Y_i$ est dans $C$, pour \unsure{tt} $i$. Mais la condition envisagée \unsure{en|dans} 2.1. \unsure{$1^\circ$|$2^\circ$}) est \deleted{\unsure{que}}\unsure{seulement} que les $u_i$ sont dans $C$. Donc en prenant pour $Y_\bullet$ un syst. inductif constant de valeur $Y$, avec $Y$ \unsure{non} $C$-cofibrant i.e. $\varnothing_M\longrightarrow Y$ non dans $C$, on a un contrexemple.

\unsure{Avec|Pour} le $=$ $I$, on voit que $u_\bullet:X_\bullet\longrightarrow Y_\bullet$ est dans $C$ ssi a) les $u_i\in C$, et b) les $u_{i,i+1}:Y_{i,i+1}\longrightarrow Y_{i+1}$ sont dans $C$. Si de plus les $u_i$ sont dans $\unsure{w|W}$, donc dans $TC$, on voit comme ci-dessus que $u_{i,i+1}\in TC_{\unsure{\varnothing}}$ pour tt $i$, donc $u_\bullet\in(TC)_I$. Donc ici on a encore

\[
 (TC)_I=C_I\cap W_{I^{\unsure{o|\circ}}},
\]

quoique $I$ ne soit pas artinien. J'ignore si l'hypothèse artinienne dans 2.1. $3^\circ$ est superflue ou non.

2) Soit $u_\bullet\in\operatorname{Fl}M(I^o)$, donc $u_\bullet$ est un objet de

\[
 \text{\uline{Fl}}\,M(I^o)\simeq M(I^o\times\unsure{\mathbb{\Delta}|\Delta}^{1})=M((I\times\mathbb{\Delta}^{1\,\mathrm{opp}})^{\mathrm{opp}}).
\]

\deleted{\unsure{P}} On vérifie que pour que cet objet soit $C$-cofibrant, il f. et s. que $1^\circ)$ $X_\bullet$ soit $C$-cofibrant, et $2^\circ)$ $u_\bullet\in C_I$ (ce qui implique que $Y_\bullet$ est aussi cofibrant)

% slanted
\margin{$I^o=I^{\mathrm{opp}}$}

% archive sheet number: 50

